\section{Anforderungen}
\begin{enumerate}
    \item Besteck kann nur von einem Philosophen gleichzeitig benutzt werden.
    \item Es gibt genauso viele Bestecke wie Philosophen.
    \item Philosophen müssen einzelne Container sein.
    \item Philosophen werden können entweder Essen oder Denken: 
    \begin{enumerate}
        \item wenn sie Denken interagieren sie nicht mit anderen Philosophen oder Besteck.
        \item wenn sie Essen wollen benötigen sie die 2 angrenzenden Bestecke um diese Tätigkeit zu vollführen und legen diese anschließend wieder ab.
    \end{enumerate}
    \item Philosophen sollen ausschließlich die 2 angrenzenden Bestecke benutzen können.
    \item Philosophen können nicht wissen wann die anderen Philosophen Essen oder Denken wollen.
    \item Philosophen müssen im Wechsel Essen und Denken können sodass sie nicht verhungern.
    \item Philosophen verhungern wenn sie nicht Essen können.
    \item Philosophen dürfen ausschließlich per Nachrichten kommunizieren und ihre Zugriffe sollen verteilt koordiniert werden.
    \item Die Performancemessungen sollen mit 5, 10 sowie einer zu ermittelnden maximal Anzahl stattfinden.
    \item Zur verteilten Koordination soll eine eigene Lösung für RPC genutzt werden und keine Bibliothek.
    \item Die eigene RPC Lösung muss mit Tests abgesichert werden.
\end{enumerate}


\subsection{Offene Fragen}
- wie sollen die Zeitspannen in den Zuständen aussehen: feste Zeitspannen, zufällig?
- Soll verhungern eine Zeitgrenze haben oder erst wenn ein deadlock entsteht?